% ====================================================================
% Supplementary Figures — per-specimen selectivity optimization
%
% One full-page figure per histology cohort sample, each showing all
% target-cluster seeds for that nerve (cross-section + optimized cuff
% pattern), ordered by achievable SI.  Generated by
% experiments_v2/figures_supp_activation.py into figures/supp/.
% ====================================================================
\clearpage
\setcounter{figure}{0}
\renewcommand{\thefigure}{S\arabic{figure}}
\section*{Supplementary Figures}
\addcontentsline{toc}{section}{Supplementary Figures}

\begin{figure}[h]
  \centering
  \figmaybe[width=0.98\textwidth]{main/fig2_phenomena.png}
  \caption{jaxon reproduces three canonical propagation phenomena
    (solid) vs.\ PyFibers/NEURON (dashed), for a myelinated A-fiber
    (MRG, top row \textbf{a}--\textbf{c}) and an unmyelinated C-fiber
    (Sundt, bottom row \textbf{d}--\textbf{f}).
    \textbf{(a, d)} AP propagation: waterfall of the membrane-potential
    trace at each node (y = position along fiber, x = time); the
    horizontal arrow marks the central stimulation site.
    \textbf{(b, e)} kHz conduction block at $20\,\mathrm{kHz}$
    extracellular hold; distal-node $V_m$ traces at four block
    amplitudes show the pace train passing, partially and fully blocking
    as amplitude increases (kHz-on window shaded).  The block threshold
    matches NEURON within the $0.01\,\mathrm{mA}$ scan resolution and the
    AP count agrees exactly.
    \textbf{(c, f)} AP collision; snapshots of $V_m$ along the fiber at
    five times (columns) for four diameters (rows) show two
    inward-propagating APs meeting and annihilating, with a
    jaxon-vs-NEURON peak-voltage discrepancy below $1\,\mu\mathrm{V}$.}
  \label{fig:phenomena}
\end{figure}

% >>> ANATOMY SECTION (auto-generated) >>>
\clearpage
\section*{Supplementary anatomy and FEM fields}
\addcontentsline{toc}{section}{Supplementary anatomy and FEM fields}

\noindent FE-model anatomy and computed extracellular fields for every nerve in the microCT cohort (Section~\ref{sec:results-duke}); the shared panel layout is described once below and applies to every page.

\newcommand{\anatcaption}{\textbf{(a)} histology slice with manual segmentation; \textbf{(b)} area-preserving deformed cross-section used to build the 3-D model; \textbf{(c)} 3-D nerve and 12-contact cuff (component legend inset); \textbf{(d)} per-contact extracellular potential $V_e$; \textbf{(e)} per-contact activating function; \textbf{(f)} per-contact extracellular E-field magnitude.}

\begin{figure}[p]\centering
  \figmaybe[width=\textwidth]{supp/figA_human_sub-47_sam-2.png}
  \caption{Specimen \texttt{human sub-47\_sam-2}: anatomy and FEM fields. \anatcaption}
  \label{fig:anat-human-sub-47-sam-2}
\end{figure}
\clearpage

\begin{figure}[p]\centering
  \figmaybe[width=\textwidth]{supp/figA_human_sub-50_sam-2.png}
  \caption{Specimen \texttt{human sub-50\_sam-2}: anatomy and FEM fields. \anatcaption}
  \label{fig:anat-human-sub-50-sam-2}
\end{figure}
\clearpage

\begin{figure}[p]\centering
  \figmaybe[width=\textwidth]{supp/figA_human_sub-53_sam-2.png}
  \caption{Specimen \texttt{human sub-53\_sam-2}: anatomy and FEM fields. \anatcaption}
  \label{fig:anat-human-sub-53-sam-2}
\end{figure}
\clearpage

\begin{figure}[p]\centering
  \figmaybe[width=\textwidth]{supp/figA_human_sub-54_sam-2.png}
  \caption{Specimen \texttt{human sub-54\_sam-2}: anatomy and FEM fields. \anatcaption}
  \label{fig:anat-human-sub-54-sam-2}
\end{figure}
\clearpage

\begin{figure}[p]\centering
  \figmaybe[width=\textwidth]{supp/figA_human_sub-54_sam-3.png}
  \caption{Specimen \texttt{human sub-54\_sam-3}: anatomy and FEM fields. \anatcaption}
  \label{fig:anat-human-sub-54-sam-3}
\end{figure}
\clearpage

\begin{figure}[p]\centering
  \figmaybe[width=\textwidth]{supp/figA_human_sub-56_sam-1.png}
  \caption{Specimen \texttt{human sub-56\_sam-1}: anatomy and FEM fields. \anatcaption}
  \label{fig:anat-human-sub-56-sam-1}
\end{figure}
\clearpage

\begin{figure}[p]\centering
  \figmaybe[width=\textwidth]{supp/figA_human_sub-56_sam-3.png}
  \caption{Specimen \texttt{human sub-56\_sam-3}: anatomy and FEM fields. \anatcaption}
  \label{fig:anat-human-sub-56-sam-3}
\end{figure}
\clearpage

\begin{figure}[p]\centering
  \figmaybe[width=\textwidth]{supp/figA_human_sub-57_sam-1.png}
  \caption{Specimen \texttt{human sub-57\_sam-1}: anatomy and FEM fields. \anatcaption}
  \label{fig:anat-human-sub-57-sam-1}
\end{figure}
\clearpage

\begin{figure}[p]\centering
  \figmaybe[width=\textwidth]{supp/figA_human_sub-57_sam-3.png}
  \caption{Specimen \texttt{human sub-57\_sam-3}: anatomy and FEM fields. \anatcaption}
  \label{fig:anat-human-sub-57-sam-3}
\end{figure}
\clearpage

\begin{figure}[p]\centering
  \figmaybe[width=\textwidth]{supp/figA_human_sub-58_sam-1.png}
  \caption{Specimen \texttt{human sub-58\_sam-1}: anatomy and FEM fields. \anatcaption}
  \label{fig:anat-human-sub-58-sam-1}
\end{figure}
\clearpage

\begin{figure}[p]\centering
  \figmaybe[width=\textwidth]{supp/figA_human_sub-58_sam-3.png}
  \caption{Specimen \texttt{human sub-58\_sam-3}: anatomy and FEM fields. \anatcaption}
  \label{fig:anat-human-sub-58-sam-3}
\end{figure}
\clearpage

\begin{figure}[p]\centering
  \figmaybe[width=\textwidth]{supp/figA_sub-10_sam-1.png}
  \caption{Specimen \texttt{swine sub-10\_sam-1}: anatomy and FEM fields. \anatcaption}
  \label{fig:anat-sub-10-sam-1}
\end{figure}
\clearpage

\begin{figure}[p]\centering
  \figmaybe[width=\textwidth]{supp/figA_sub-11_sam-1.png}
  \caption{Specimen \texttt{swine sub-11\_sam-1}: anatomy and FEM fields. \anatcaption}
  \label{fig:anat-sub-11-sam-1}
\end{figure}
\clearpage

\begin{figure}[p]\centering
  \figmaybe[width=\textwidth]{supp/figA_sub-11_sam-3.png}
  \caption{Specimen \texttt{swine sub-11\_sam-3}: anatomy and FEM fields. \anatcaption}
  \label{fig:anat-sub-11-sam-3}
\end{figure}
\clearpage

\begin{figure}[p]\centering
  \figmaybe[width=\textwidth]{supp/figA_sub-12_sam-1.png}
  \caption{Specimen \texttt{swine sub-12\_sam-1}: anatomy and FEM fields. \anatcaption}
  \label{fig:anat-sub-12-sam-1}
\end{figure}
\clearpage

\begin{figure}[p]\centering
  \figmaybe[width=\textwidth]{supp/figA_sub-12_sam-3.png}
  \caption{Specimen \texttt{swine sub-12\_sam-3}: anatomy and FEM fields. \anatcaption}
  \label{fig:anat-sub-12-sam-3}
\end{figure}
\clearpage

\begin{figure}[p]\centering
  \figmaybe[width=\textwidth]{supp/figA_sub-13_sam-1.png}
  \caption{Specimen \texttt{swine sub-13\_sam-1}: anatomy and FEM fields. \anatcaption}
  \label{fig:anat-sub-13-sam-1}
\end{figure}
\clearpage

\begin{figure}[p]\centering
  \figmaybe[width=\textwidth]{supp/figA_sub-13_sam-3.png}
  \caption{Specimen \texttt{swine sub-13\_sam-3}: anatomy and FEM fields. \anatcaption}
  \label{fig:anat-sub-13-sam-3}
\end{figure}
\clearpage

\begin{figure}[p]\centering
  \figmaybe[width=\textwidth]{supp/figA_sub-14_sam-2.png}
  \caption{Specimen \texttt{swine sub-14\_sam-2}: anatomy and FEM fields. \anatcaption}
  \label{fig:anat-sub-14-sam-2}
\end{figure}
\clearpage

\begin{figure}[p]\centering
  \figmaybe[width=\textwidth]{supp/figA_sub-14_sam-3.png}
  \caption{Specimen \texttt{swine sub-14\_sam-3}: anatomy and FEM fields. \anatcaption}
  \label{fig:anat-sub-14-sam-3}
\end{figure}
\clearpage

\begin{figure}[p]\centering
  \figmaybe[width=\textwidth]{supp/figA_sub-15_sam-2.png}
  \caption{Specimen \texttt{swine sub-15\_sam-2}: anatomy and FEM fields. \anatcaption}
  \label{fig:anat-sub-15-sam-2}
\end{figure}
\clearpage

\begin{figure}[p]\centering
  \figmaybe[width=\textwidth]{supp/figA_sub-15_sam-3.png}
  \caption{Specimen \texttt{swine sub-15\_sam-3}: anatomy and FEM fields. \anatcaption}
  \label{fig:anat-sub-15-sam-3}
\end{figure}
\clearpage

\begin{figure}[p]\centering
  \figmaybe[width=\textwidth]{supp/figA_sub-4_sam-3.png}
  \caption{Specimen \texttt{swine sub-4\_sam-3}: anatomy and FEM fields. \anatcaption}
  \label{fig:anat-sub-4-sam-3}
\end{figure}
\clearpage

\begin{figure}[p]\centering
  \figmaybe[width=\textwidth]{supp/figA_sub-5_sam-1.png}
  \caption{Specimen \texttt{swine sub-5\_sam-1}: anatomy and FEM fields. \anatcaption}
  \label{fig:anat-sub-5-sam-1}
\end{figure}
\clearpage

\begin{figure}[p]\centering
  \figmaybe[width=\textwidth]{supp/figA_sub-5_sam-3.png}
  \caption{Specimen \texttt{swine sub-5\_sam-3}: anatomy and FEM fields. \anatcaption}
  \label{fig:anat-sub-5-sam-3}
\end{figure}
\clearpage

\begin{figure}[p]\centering
  \figmaybe[width=\textwidth]{supp/figA_sub-6_sam-7.png}
  \caption{Specimen \texttt{swine sub-6\_sam-7}: anatomy and FEM fields. \anatcaption}
  \label{fig:anat-sub-6-sam-7}
\end{figure}
\clearpage

\begin{figure}[p]\centering
  \figmaybe[width=\textwidth]{supp/figA_sub-8_sam-1.png}
  \caption{Specimen \texttt{swine sub-8\_sam-1}: anatomy and FEM fields. \anatcaption}
  \label{fig:anat-sub-8-sam-1}
\end{figure}
\clearpage

\begin{figure}[p]\centering
  \figmaybe[width=\textwidth]{supp/figA_sub-8_sam-7.png}
  \caption{Specimen \texttt{swine sub-8\_sam-7}: anatomy and FEM fields. \anatcaption}
  \label{fig:anat-sub-8-sam-7}
\end{figure}
\clearpage

\begin{figure}[p]\centering
  \figmaybe[width=\textwidth]{supp/figA_sub-9_sam-3.png}
  \caption{Specimen \texttt{swine sub-9\_sam-3}: anatomy and FEM fields. \anatcaption}
  \label{fig:anat-sub-9-sam-3}
\end{figure}
\clearpage
% <<< ANATOMY SECTION <<<


\noindent Per-specimen selectivity-optimization outcomes for the full
histology-derived vagus-nerve cohort
(Section~\ref{sec:results-duke}).  Each figure collects all
target-cluster seeds for one nerve; the shared layout and color
scheme are described once below and apply to every panel.

% Shared caption text (no arguments -> no catcode issues with the
% underscores in sample identifiers).
\newcommand{\suppcaption}{All target-cluster seeds for this nerve,
  ordered left-to-right by achievable selectivity index $\mathrm{SI}$.
  Each seed shows the nerve cross-section (top) with fibers colored by
  post-optimization recruitment (target fired / target silent /
  off-target fired / off-target silent) and the peripheral
  target-fascicle cluster outlined, above the optimized 12-contact
  cuff pattern (bottom; cathode blue, anode orange, inactive white,
  marker size $\propto |I_k|$, over four angular columns $\times$ three
  axial rows).}

% ── Swine subset ──────────────────────────────────────────────────────
\begin{figure}[p]\centering
  \figmaybe[width=0.95\textwidth]{supp/figS_sub-4_sam-3.png}
  \caption{Specimen \texttt{swine sub-4\_sam-3}. \suppcaption}
  \label{fig:supp-sub-4-sam-3}
\end{figure}

\begin{figure}[p]\centering
  \figmaybe[width=0.95\textwidth]{supp/figS_sub-5_sam-1.png}
  \caption{Specimen \texttt{swine sub-5\_sam-1}. \suppcaption}
  \label{fig:supp-sub-5-sam-1}
\end{figure}

\begin{figure}[p]\centering
  \figmaybe[width=0.95\textwidth]{supp/figS_sub-5_sam-3.png}
  \caption{Specimen \texttt{swine sub-5\_sam-3}. \suppcaption}
  \label{fig:supp-sub-5-sam-3}
\end{figure}

\begin{figure}[p]\centering
  \figmaybe[width=0.95\textwidth]{supp/figS_sub-6_sam-7.png}
  \caption{Specimen \texttt{swine sub-6\_sam-7}. \suppcaption}
  \label{fig:supp-sub-6-sam-7}
\end{figure}

\begin{figure}[p]\centering
  \figmaybe[width=0.95\textwidth]{supp/figS_sub-8_sam-1.png}
  \caption{Specimen \texttt{swine sub-8\_sam-1}. \suppcaption}
  \label{fig:supp-sub-8-sam-1}
\end{figure}

\begin{figure}[p]\centering
  \figmaybe[width=0.95\textwidth]{supp/figS_sub-8_sam-7.png}
  \caption{Specimen \texttt{swine sub-8\_sam-7}. \suppcaption}
  \label{fig:supp-sub-8-sam-7}
\end{figure}

\begin{figure}[p]\centering
  \figmaybe[width=0.95\textwidth]{supp/figS_sub-9_sam-3.png}
  \caption{Specimen \texttt{swine sub-9\_sam-3}. \suppcaption}
  \label{fig:supp-sub-9-sam-3}
\end{figure}

\begin{figure}[p]\centering
  \figmaybe[width=0.95\textwidth]{supp/figS_sub-10_sam-1.png}
  \caption{Specimen \texttt{swine sub-10\_sam-1}. \suppcaption}
  \label{fig:supp-sub-10-sam-1}
\end{figure}

\begin{figure}[p]\centering
  \figmaybe[width=0.95\textwidth]{supp/figS_sub-11_sam-1.png}
  \caption{Specimen \texttt{swine sub-11\_sam-1}. \suppcaption}
  \label{fig:supp-sub-11-sam-1}
\end{figure}

\begin{figure}[p]\centering
  \figmaybe[width=0.95\textwidth]{supp/figS_sub-11_sam-3.png}
  \caption{Specimen \texttt{swine sub-11\_sam-3}. \suppcaption}
  \label{fig:supp-sub-11-sam-3}
\end{figure}

\begin{figure}[p]\centering
  \figmaybe[width=0.95\textwidth]{supp/figS_sub-12_sam-1.png}
  \caption{Specimen \texttt{swine sub-12\_sam-1}. \suppcaption}
  \label{fig:supp-sub-12-sam-1}
\end{figure}

\begin{figure}[p]\centering
  \figmaybe[width=0.95\textwidth]{supp/figS_sub-12_sam-3.png}
  \caption{Specimen \texttt{swine sub-12\_sam-3}. \suppcaption}
  \label{fig:supp-sub-12-sam-3}
\end{figure}

\begin{figure}[p]\centering
  \figmaybe[width=0.95\textwidth]{supp/figS_sub-13_sam-1.png}
  \caption{Specimen \texttt{swine sub-13\_sam-1}. \suppcaption}
  \label{fig:supp-sub-13-sam-1}
\end{figure}

\begin{figure}[p]\centering
  \figmaybe[width=0.95\textwidth]{supp/figS_sub-13_sam-3.png}
  \caption{Specimen \texttt{swine sub-13\_sam-3}. \suppcaption}
  \label{fig:supp-sub-13-sam-3}
\end{figure}

\begin{figure}[p]\centering
  \figmaybe[width=0.95\textwidth]{supp/figS_sub-14_sam-2.png}
  \caption{Specimen \texttt{swine sub-14\_sam-2}. \suppcaption}
  \label{fig:supp-sub-14-sam-2}
\end{figure}

\begin{figure}[p]\centering
  \figmaybe[width=0.95\textwidth]{supp/figS_sub-14_sam-3.png}
  \caption{Specimen \texttt{swine sub-14\_sam-3}. \suppcaption}
  \label{fig:supp-sub-14-sam-3}
\end{figure}

\begin{figure}[p]\centering
  \figmaybe[width=0.95\textwidth]{supp/figS_sub-15_sam-2.png}
  \caption{Specimen \texttt{swine sub-15\_sam-2}. \suppcaption}
  \label{fig:supp-sub-15-sam-2}
\end{figure}

\begin{figure}[p]\centering
  \figmaybe[width=0.95\textwidth]{supp/figS_sub-15_sam-3.png}
  \caption{Specimen \texttt{swine sub-15\_sam-3}. \suppcaption}
  \label{fig:supp-sub-15-sam-3}
\end{figure}

% ── Human subset ──────────────────────────────────────────────────────
\begin{figure}[p]\centering
  \figmaybe[width=0.95\textwidth]{supp/figS_human_sub-47_sam-2.png}
  \caption{Specimen \texttt{human sub-47\_sam-2}. \suppcaption}
  \label{fig:supp-human-sub-47-sam-2}
\end{figure}

\begin{figure}[p]\centering
  \figmaybe[width=0.95\textwidth]{supp/figS_human_sub-50_sam-2.png}
  \caption{Specimen \texttt{human sub-50\_sam-2}. \suppcaption}
  \label{fig:supp-human-sub-50-sam-2}
\end{figure}

\begin{figure}[p]\centering
  \figmaybe[width=0.95\textwidth]{supp/figS_human_sub-53_sam-2.png}
  \caption{Specimen \texttt{human sub-53\_sam-2}. \suppcaption}
  \label{fig:supp-human-sub-53-sam-2}
\end{figure}

\begin{figure}[p]\centering
  \figmaybe[width=0.95\textwidth]{supp/figS_human_sub-54_sam-2.png}
  \caption{Specimen \texttt{human sub-54\_sam-2}. \suppcaption}
  \label{fig:supp-human-sub-54-sam-2}
\end{figure}

\begin{figure}[p]\centering
  \figmaybe[width=0.95\textwidth]{supp/figS_human_sub-54_sam-3.png}
  \caption{Specimen \texttt{human sub-54\_sam-3}. \suppcaption}
  \label{fig:supp-human-sub-54-sam-3}
\end{figure}

\begin{figure}[p]\centering
  \figmaybe[width=0.95\textwidth]{supp/figS_human_sub-56_sam-1.png}
  \caption{Specimen \texttt{human sub-56\_sam-1}. \suppcaption}
  \label{fig:supp-human-sub-56-sam-1}
\end{figure}

\begin{figure}[p]\centering
  \figmaybe[width=0.95\textwidth]{supp/figS_human_sub-56_sam-3.png}
  \caption{Specimen \texttt{human sub-56\_sam-3}. \suppcaption}
  \label{fig:supp-human-sub-56-sam-3}
\end{figure}

\begin{figure}[p]\centering
  \figmaybe[width=0.95\textwidth]{supp/figS_human_sub-57_sam-3.png}
  \caption{Specimen \texttt{human sub-57\_sam-3}. \suppcaption}
  \label{fig:supp-human-sub-57-sam-3}
\end{figure}

\begin{figure}[p]\centering
  \figmaybe[width=0.95\textwidth]{supp/figS_human_sub-58_sam-3.png}
  \caption{Specimen \texttt{human sub-58\_sam-3}. \suppcaption}
  \label{fig:supp-human-sub-58-sam-3}
\end{figure}
